\section{Conclusion and Future Work}
CloudCanvas implements an AWS-focused workflow that connects an editable resource graph, structured AI proposals, explicit application and publish decisions, persisted lifecycle records, and dependency-aware SDK operations. The integration is the system contribution; visual cloud modeling, topology execution, and AI-assisted infrastructure generation remain established prior work. The paper's experiment examines the shared graph contract rather than attributing broad correctness or productivity gains to the interface.

For RQ1, the full contract correctly classifies all 72 synthetic conformance instances, compared with $55.56\%$ for schema-only validation and $77.78\%$ for the schema-plus-topology ablation, with no observed false rejections. For RQ2, full validation costs $72.9\pm14.1$ microseconds per small graph across three fresh processes, while all six selected boundary controls remain accepted, demonstrating that the structural predicate has limited operational meaning. These numbers describe the supplied fixtures and workstation; they do not establish superiority in live deployment, security, intent satisfaction, or usability.

Three concrete directions follow. First, add independent graph-intent and policy oracles, then evaluate catalog grounding and review with matched models and held-out tasks. Second, make execution durable through atomic operation reservation, idempotency, complete failure finalization, and interrupted-operation recovery. Third, conduct bounded live-AWS lifecycle and controlled user studies that report readiness, cleanup, unsafe approvals, and correction effort separately. Those studies would test whether the implemented workflow produces benefits beyond the conformance boundary demonstrated here.
